\documentclass[11pt]{article}
\usepackage[T1]{fontenc}
\usepackage{amsmath,amssymb,amsthm}
\usepackage[margin=1in]{geometry}
\usepackage{xurl}
\usepackage[hidelinks]{hyperref}
\newtheorem{theorem}{Theorem}
\newtheorem{corollary}[theorem]{Corollary}
\newtheorem{lemma}[theorem]{Lemma}
\title{A classical segment for the local--global passage\\
\large From primewise neighbourhoods to a finite \'etale cover}
\author{}
\date{31 October 2026}
\begin{document}
\maketitle

\part{The mathematical argument}

\section{The precise classical principle}

We work with commutative rings with identity, including the zero ring.
In this part we assume classical logic and the following principle:
\begin{quote}
Every proper ideal of a ring is contained in a prime ideal.
\end{quote}
The principle is kept separate from classical logic. It can subsequently
be deduced from choice by applying Zorn's lemma to the proper ideals
containing a given proper ideal, and proving that a maximal ideal is
prime. That derivation is not needed for the argument below.

\begin{theorem}\label{thm:cover}
Let $E$ be a ring and let $G$ be a property of its elements. Suppose
that for every prime ideal $\mathfrak q$ there exists $s\in E$ such that
$G(s)$ and $s\notin\mathfrak q$. Then there are finitely many elements
$s_1,\ldots,s_m$, each satisfying $G$, and coefficients $a_i\in E$ with
\[
 \sum_{i=1}^m s_i a_i=1.
\]
The family may be empty when $E$ is the zero ring.
\end{theorem}
\begin{proof}
Let $J$ be the ideal generated by all elements satisfying $G$. Membership
in $J$ means having a finite expression $\sum_i s_i a_i$ with $G(s_i)$.
If $1\notin J$, prime-ideal existence supplies $\mathfrak q\supseteq J$.
The hypothesis gives a good $s$ outside $\mathfrak q$, while $s\in J$
puts it inside $\mathfrak q$, a contradiction. Classical logic therefore
gives $1\in J$. Its finite expression is exactly the required equation.
\end{proof}

This proof never chooses a neighbourhood simultaneously for every prime.
It considers just one prime containing $J$ under the assumption that $J$
is proper. The resulting equation is the algebraic certificate that
the principal opens $D(s_i)$ cover $\operatorname{Spec}E$.

The property may mean that an element carries additional data. After the
finite family is obtained, finitely many witnesses for that data can be
retained along with the equation. This uses only finite choice, proved
by induction on the family length.

\section{Application to the local structure theorem}

Fix a ring map $R\to E$. In the monic convention of
\emph{Local structure of \'etale algebras}, a standard \'etale algebra has
the form
\[
 (R[T]/(P))[1/G(t)],
\]
where $P$ is monic and $P'(t)$ is invertible in the displayed algebra.
Inverting $G(t)P'(t)$ instead gives the same kind of presentation.

\begin{corollary}\label{cor:etale}
Suppose that for every prime $\mathfrak q$ of $E$ there exists
$s\notin\mathfrak q$ such that $E[1/s]$ is standard \'etale over $R$.
Then there are $s_1,\ldots,s_m$ generating the unit ideal of $E$ such
that every $E[1/s_i]$ is standard \'etale over $R$.
The standard presentations and localisation isomorphisms may be
retained as part of the conclusion.
\end{corollary}
\begin{proof}
Apply Theorem~\ref{thm:cover} to the property that $E[1/s]$ has such a
presentation. Retain its finite family of presentation witnesses.
\end{proof}

There is a distinction between a neighbourhood standard over $R$ and
one standard only over $R_{\mathfrak p}$. To apply the corollary to the
local theorem used in the ZMT argument, one first carries out the
following algebraic steps.

Given $\mathfrak q\subset E$, put $\mathfrak p=\mathfrak q\cap R$.
Classically $R_{\mathfrak p}$ is residually discrete local. The local
theorem, after its hypotheses have been supplied, gives a principal
neighbourhood standard over $R_{\mathfrak p}$. Clear the denominator
of its defining element to obtain $s\in E\setminus\mathfrak q$.

Finite presentation is then used to descend the presentation and its
inverse maps: all the data and identities involve finitely many
coefficients. Consequently there is $f\in R\setminus\mathfrak p$ such
that $E[1/sf]$ is standard \'etale over $R[1/f]$.
The equations certifying the inverse maps and the inverse derivative
must be included among the identities being descended.

Finally change the base back to $R$ by monic rescaling. If the equation
over $R[1/f]$ is $P(T)$ of degree $d>0$, choose $N$ large enough that
\[
 \widetilde P(Z)=f^{Nd}P(Z/f^N)
\]
has a monic lift to $R[Z]$. If $G$ is the polynomial being inverted,
choose $M$ such that $f^M G(Z/f^N)$ has a lift $V\in R[Z]$. Then
\[
 (R[1/f][T]/(P))[1/G(t)]
 \simeq (R[Z]/(\widetilde P))[1/(fV(z))].
\]
The factor $f$ in the denominator is retained. It makes the substitution
$z=f^Nt$ invertible and removes any ambiguity from the chosen lifts.
The derivative becomes $f^{N(d-1)}P'(t)$, so remains invertible.
The degree-zero case is the zero algebra and is treated separately.

This supplies the hypothesis of Corollary~\ref{cor:etale}. Thus classical
prime-ideal existence replaces the dynamical local--global argument;
finite-presentation descent and the change-of-base lemma remain parts
of the proof.

\part{Connection with the formalisation}

\section{The classical segment and its parameters}

The development is in \path{unit/sag/etale-classical-global}.
The existing constructive imports remain in the usual \texttt{Theory}
segment, with univalence, replacement and surjection parameters.
Inside it is the separate segment
\begin{center}
\texttt{Classical.Cl[lem,prime-existence]}.
\end{center}
Only imports within this segment can use these two parameters.
Neither principle is added as a global axiom or to the constructive
ZMT or local \'etale developments.

\texttt{lem} is excluded middle for propositions at each universe level.
\texttt{prime-existence} asserts the mere existence of a prime above
each proper proposition-valued ideal of each commutative set-ring.
The type of this principle is defined in \texttt{PrimeExistence},
outside the classical segment. No inhabitant is asserted there.

A prime in this interface has propositional membership, properness and
mere disjunctive primality. Decidability is not part of the existence
principle. \texttt{ClassicalLogic.Prime} derives it from excluded middle
and constructs the exact \texttt{PrimeIdempotents.Prime} record used by
the earlier local proofs. Deriving prime-ideal existence from choice is
left as a future theorem.

\section{The checked compactness argument}

The generic theorem is
\begin{center}
\texttt{ClassicalPrimeCover.Over[E,setE,G]}\\
\texttt{.Given[pointwise].theorem}.
\end{center}
Its premise is the truncated existence of a good element outside each
prime. Its conclusion is the truncated existence of a natural number
$m$, two families $s,a:\mathrm{Fin}(m)\to E$, witnesses of $G(s_i)$,
and the equation $\sum_i s_i a_i=1$.

The ideal of good elements is constructed using finite lists and
propositional truncation. Its ideal laws and the inclusion of every good
element are proved. A hypothetical proper ideal gives a prime above it;
the primewise hypothesis contradicts containment. Double-negation
elimination is applied only to the proposition $1\in J$. The finite
certificate is then extracted under truncation. Compactness itself is
proved, rather than introduced as a further parameter.

Actual localisation certificates contain algebras and universal
properties, so live one universe above the ring carrier. The variant
\texttt{ClassicalCertifiedCover} handles this without a resizing axiom.
For each $s$, excluded middle represents the proposition
$\|\mathrm{Certificate}(s)\|$ by either the small type \texttt{Top}
or the small type \texttt{Empty}, with maps in both directions. These
small predicates generate the ideal. Once the finite family is found,
ordinary finite choice restores its actual certificates inside the
final truncation. This step uses no additional axiom of choice.

\section{The \'etale endpoints and their scope}

\texttt{ClassicalEtaleCover.Over[A,setA,E,setE,iota]} exports three
adapters. Each takes a neighbourhood theorem for every decidable prime
of $E$ and gives a finite cover with its coefficient equation.
\begin{itemize}
\item \texttt{Given[local].theorem} accepts the exact conclusion of
\texttt{Local.Statement}. Every chart includes an actual algebra $T$,
the map $E\to T$, the universal property $T=E[1/s]$, and the existing
\texttt{StandardEtale.Points.Standard} certificate.
\item \texttt{GivenFlat[local].theorem} accepts the flat local conclusion.
It additionally retains the square-zero lifting definition of \'etale
for each chart.
\item \texttt{GivenMonic[local].theorem} accepts and retains a stronger
\texttt{MonicChart}: its polynomial is explicitly a coefficient list
followed by the leading coefficient $1$. This supplies the monic form
of Corollary~\ref{cor:etale}.
\end{itemize}

The last distinction matters because the older \texttt{Local.Statement}
conclusion quantifies its polynomial $Q$ but does not record monicity.
The new ordinary adapter preserves that interface; it does not recover
missing information. A forgetful map from \texttt{MonicChart} to the
older chart interface is proved.

These are conditional local-to-global theorems: their neighbourhoods
must already be standard over the original base $A$. This development
does not construct the local setup at $A_{\mathfrak p}$, discharge the
older local proof's remaining inputs, or formalise the finite-character
and monic rescaling lemmas described in Part~I. It also does not assemble
a global square-zero lifting certificate from the finite cover.
Accordingly it proves the classical compactness step, while the full
global structure theorem from its original hypotheses still needs that
integration. The constructive dynamical replacement remains separate.

\section{Validation}

A fresh check of \texttt{All.seg} gives $4262$ declarations, no failures,
in $1.49$ seconds. The opening audit covers all $217$ named declarations
in \texttt{PrimeExistence} and \texttt{Classical}. Each opens at most
$11$ definitions and is reported at $0.00$ seconds with two-decimal
timing. The audit reports no openings outside the specified lists.

Regressions cover the good element $1$, the empty cover of the zero ring,
the impossibility of an empty-good cover in a nonzero ring, pairs
$a+b=1$ and complementary elements $e,1-e$, and certificates containing
a universe-valued component. An actual standard algebra for $Q=X$,
$G=1$ is constructed over every base using \texttt{StandardAlgebra.Build}.
Its identity localisation at $1$ is proved by the universal property,
and both the ordinary and monic cover theorems are applied to it.

Three mutations in isolated copies were rejected: changing the cover
equation to a sum equal to zero, applying prime properness to the wrong
element, and replacing the leading coefficient $1$ by $0$. The last is
caught by the constructed-standard regression. Shared sources and the
checker are unchanged; reused sources are relative links.

From \path{unit/}, check the development with
\begin{verbatim}
./segc sag/etale-classical-global all
./segc sag/etale-classical-global stop
\end{verbatim}
The exact signatures and the fresh-check command are in
\path{unit/sag/etale-classical-global/README.txt}.

\end{document}
